\documentclass[11pt]{article}
\usepackage[T1]{fontenc}
\usepackage[margin=1in]{geometry}
\usepackage{amsmath,amssymb}
\usepackage[hidelinks]{hyperref}

\newcommand{\SA}{\operatorname{SA}}
\newcommand{\Lip}{\operatorname{Lip}}
\newcommand{\co}{\operatorname{co}}
\newcommand{\strexp}{\operatorname{str\mbox{-}exp}}
\newcommand{\F}{\mathcal F}
\newcommand{\R}{\mathbb R}

\title{Literature Status: Strongly Norm Attaining Lipschitz Maps}
\author{Automatic Math Discovery Run \texttt{fa\_banach\_001}}
\date{June 16, 2026}

\begin{document}
\maketitle

\section*{Status}

This is a \textbf{literature\_already\_answered} packet.  It records that the
main open-problem signals extracted from arXiv:1807.03363 were explicitly
answered or sharply delimited by the later paper arXiv:1907.07698.  It is not
a new proof packet.

\section*{Source Paper and Questions}

The source paper is:

\begin{quote}
B. Cascales, R. Chiclana, L. C. Garcia-Lirola, M. Martin, and
A. Rueda Zoca,
\emph{On strongly norm attaining Lipschitz maps}, arXiv:1807.03363.
\end{quote}

In the introduction, the authors note that if
\(\SA(M,Y)\) is dense in \(\Lip(M,Y)\) for every Banach space \(Y\), then
\(\F(M)\) has Lindenstrauss property A.  They state that the converse is not
known and that, as far as they know, the Radon--Nikodym property of
\(\F(M)\) could be necessary for this density.  This is the source-paper signal
covered by the present literature-status packet.

\section*{Supporting Answer Paper}

The supporting paper is:

\begin{quote}
R. Chiclana, L. C. Garcia-Lirola, M. Martin, A. Rueda Zoca, and
P. Tradacete,
\emph{Examples and applications of the density of strongly norm attaining
Lipschitz maps}, arXiv:1907.07698.
\end{quote}

The introduction of arXiv:1907.07698 formulates the relevant issues as
Questions Q1--Q4.  In particular, Q1 asks whether \(\F(M)\) must have the RNP
when \(\SA(M,Y)\) is dense in \(\Lip(M,Y)\) for every Banach space \(Y\);
Q2 asks whether
\[
  B_{\F(M)}
  =
  \overline{\co}\bigl(\strexp(B_{\F(M)})\bigr)
\]
or Gromov concavity suffices for density of scalar strongly norm attaining
Lipschitz maps; Q3 asks for the converse implication; and Q4 asks whether the
presence of an isomorphic, or even isometric, \(L_1[0,1]\) copy forces failure
of density.

\section*{Answer Identification}

Theorem \texttt{theo:lluvia} of arXiv:1907.07698 constructs metric spaces
\(M_p\), \(p=1,2\), such that \(\SA(M_p,Y)\) is dense in \(\Lip(M_p,Y)\) for
every Banach space \(Y\), while each \(M_p\) contains an isometric copy of
\([0,1]\).  Consequently \(\F(M_p)\) contains a copy of \(L_1[0,1]\) and fails
the RNP.  This gives a negative answer to the source paper's RNP-necessity
question, and also a negative answer to Q4.

Theorem \texttt{theo:toro} of arXiv:1907.07698 treats the Euclidean unit
circle \(\mathbb T\).  It proves that \(\SA(\mathbb T,\R)\) is not dense in
\(\Lip(\mathbb T,\R)\), although every molecule is a strongly exposed point of
\(B_{\F(\mathbb T)}\).  Thus \(B_{\F(\mathbb T)}\) is the closed convex hull of
its strongly exposed points, and \(\mathbb T\) is Gromov concave.  This gives a
negative answer to Q2.

For Q3, Theorem \texttt{theorem:clcvstrepcompact} of arXiv:1907.07698 proves
that if \(M\) is compact and \(\SA(M,\R)\) is dense in \(\Lip(M,\R)\), then
\[
  B_{\F(M)}
  =
  \overline{\co}\bigl(\strexp(B_{\F(M)})\bigr).
\]
Corollary \texttt{coro:clcvstrepbouncompact} extends this conclusion to
boundedly compact metric spaces.

\section*{Scope Limitations}

This packet clears the RNP/density branch from arXiv:1807.03363 and records
closely related consequences developed in arXiv:1907.07698.  It does not claim
to solve the newer residual questions raised inside arXiv:1907.07698, such as
the bi-Lipschitz closed-subset variants of the RNP question or the question
whether an equivalent metric on \([0,1]\) can yield density of
\(\SA(([0,1],d),\R)\).  It also does not record a later answer to the separate
octahedrality question mentioned in arXiv:1807.03363.

\section*{Search Evidence}

Local cheap indexes for arXiv:1807.03363, ``strongly norm attaining Lipschitz
maps'', ``Lipschitz-free'', and ``RNP'' had no prior exact packet at
reservation time.  Direct source inspection found arXiv:1907.07698 as a later
paper by overlapping authors that explicitly cites arXiv:1807.03363 and
formulates the relevant density/RNP question as Q1 before answering it as
described above.

\begin{thebibliography}{9}

\bibitem{CCGMR1807}
B. Cascales, R. Chiclana, L. C. Garcia-Lirola, M. Martin, and
A. Rueda Zoca,
\emph{On strongly norm attaining Lipschitz maps}, arXiv:1807.03363.

\bibitem{CGMRT1907}
R. Chiclana, L. C. Garcia-Lirola, M. Martin, A. Rueda Zoca, and
P. Tradacete,
\emph{Examples and applications of the density of strongly norm attaining
Lipschitz maps}, arXiv:1907.07698.

\end{thebibliography}

\end{document}
